% pruning_bounds.tex -- how the ttauto search decides a path is hopeless,
% and a sharper test (issue #23).
%
% Build: latexmk -pdf pruning_bounds.tex
% The tables are produced by make_tables.py from bounds_toy.cpp.
\documentclass[11pt]{amsart}
\usepackage{amsmath,amssymb,amsthm}
\usepackage[margin=1in]{geometry}
\usepackage{booktabs}
\usepackage[table]{xcolor}
\usepackage{hyperref}
\hypersetup{colorlinks=true,linkcolor=blue!60!black,citecolor=blue!60!black,
  urlcolor=blue!60!black}

\newcommand{\code}[1]{\texttt{#1}}
\newcommand{\ttauto}{\texttt{ttauto}}
\newcommand{\lam}{\lambda}
\newcommand{\Lam}{\Lambda}
\newcommand{\nrm}[1]{\lVert #1\rVert}
\newcommand{\D}{D}
\newcommand{\Pos}{\mathcal{P}}
\newcommand{\Hgp}{\mathcal{H}}
\newcommand{\tr}{\operatorname{tr}}

% Table highlighting (used by tables.tex, which make_tables.py writes).
\definecolor{goodbg}{rgb}{0.80,0.93,0.80}
\definecolor{caughtbg}{rgb}{0.99,0.87,0.70}
\newcommand{\good}[1]{\cellcolor{goodbg}\textbf{#1}}
\newcommand{\caught}[1]{\cellcolor{caughtbg}#1}

\theoremstyle{plain}
\newtheorem{Theorem}{Theorem}[section]
\newtheorem{Proposition}[Theorem]{Proposition}
\newtheorem{Lemma}[Theorem]{Lemma}
\newtheorem{Corollary}[Theorem]{Corollary}
\theoremstyle{definition}
\newtheorem{Definition}[Theorem]{Definition}
\newtheorem{Remark}[Theorem]{Remark}
\newtheorem{Example}[Theorem]{Example}

\begin{document}

\title{Pruning bounds for folding-automaton searches}
\author{Notes for \ttauto\ issue \#23}
\date{\today}

\begin{abstract}
  To list the pseudo-Anosov maps with dilatation at most~$\Lam$, \ttauto\
  explores the paths of a folding automaton and abandons a path as soon as
  it can prove that no way of continuing it back to its start leads to a
  dilatation within the window.  We set out the tests the code uses for
  this and why each is safe.  They lose not in the Ham--Song lemma itself
  but in applying it to paths that are only partly built: a run of shears
  grows the norm while the dilatation stays at one.  We then prove a
  sharper test, a lower bound on the dilatation of every possible way of
  closing up a path, and check it with a small program that knows nothing
  of \ttauto.  On the examples it removes all, or nearly all, of the
  search that leads nowhere.
\end{abstract}

\maketitle

\section{Summary}

The search spends most of its time on partial paths that can never close
up within the window.  In the three-puncture automaton, for instance, no
accepted closed path is longer than $T-1$ for traces up to~$T$, yet the
mean length of the partial paths the search explores grows roughly like
$T^2$, reaching 677 at $T=80$.  The cause is that the tests in use are all
lower bounds on quantities of the partial path itself, and along a run of
shears those stay small.

The new test (Theorems~\ref{thm:rauzy} and~\ref{thm:general}) bounds
instead the dilatation of every completion.  A completion must make the
matrix irreducible, so it must add enough entries to make the partial
path's digraph strongly connected, and that forces its dilatation up.  In
dimension two it cuts a run of shears exactly where it stops being useful
(Theorem~\ref{thm:d2}).  On the examples of Section~\ref{sec:check} it
finds the same pseudo-Anosov maps as the current tests, never missing one,
while exploring up to $10^5$ times fewer partial paths.  Its cost in the
larger automata of \ttauto\ is the main open question
(Section~\ref{sec:open}).

\section{Setting}

A folding automaton is a finite directed graph whose edges are labelled by
\emph{fold matrices}
\begin{equation}
  \label{eq:fold}
  F = P + e_{ab},
\end{equation}
where $P$ is a $d\times d$ permutation matrix, $d$ is the number of main
edges of the train tracks, and $e_{ab}$ has a single~$1$, in row~$a$ and
column~$b$, where $P$ has a zero.  (In the code these are
\code{mathmatrix\_permplus1}.)  We write $P(c)$ for the index with
$Pe_c=e_{P(c)}$.  The search multiplies on the left
(\code{TMl.push(F * TMl.top())} in \code{ttauto.hpp}), so a path with folds
$F_1,\dots,F_k$, in that order, has matrix
\[
  M = F_k F_{k-1}\cdots F_1 .
\]
A closed path is \emph{accepted} when $M$ is primitive and
$1<\rho(M)\le\Lam$, where $\rho$ is the spectral radius.  (\ttauto\ also
applies the gate test of issue~\#2, which plays no part here.)

The search is depth first from a start vertex.  A \emph{prefix} is a path
from the start vertex, with matrix~$A$.  A \emph{completion} of it is a path
from its final vertex back to the start, with matrix~$B$, possibly empty
when the prefix is already closed.  Following the prefix by a completion
gives a closed path with matrix $BA$, and every closed path that begins
with the prefix arises this way, from exactly one completion.

A \emph{test} examines a prefix and may declare it hopeless, in which case
the search abandons it and everything beyond it.  A test is \emph{safe} if
it never abandons a prefix that some completion turns into an accepted
closed path.  An unsafe test makes the search miss pseudo-Anosov maps
without saying so, so safety comes first; speed is what we then try to
gain.

\emph{Notation.}  For a nonnegative matrix $X$: $\nrm{X}$ is the sum of its
entries; $r(X)$ and $c(X)$ are its smallest row sum and smallest column
sum; $\D(X)$ is its digraph, with an edge $j\to i$ whenever $X_{ij}>0$.  A
\emph{position} is a pair $(a,b)$, the place of the $1$ in $e_{ab}$; for a
set $S$ of positions, $E_S=\sum_{(a,b)\in S}e_{ab}$, and $\D(S)$ has the
edges $b\to a$, so that $\D(E_S)=\D(S)$.  Two facts about nonnegative
matrices are used throughout~\cite{Seneta}: if $0\le X\le Y$ entrywise then
$\rho(X)\le\rho(Y)$; and for every vector $x>0$,
\begin{equation}
  \label{eq:cw}
  \min_i \frac{(Xx)_i}{x_i} \;\le\; \rho(X),
\end{equation}
the Collatz--Wielandt bound (the maximum over $i$ bounds $\rho(X)$ above
when $X$ is irreducible).  With $x=\mathbf 1$ it gives $r(X)\le\rho(X)$,
and applied to $X^{\mathsf T}$ it gives $c(X)\le\rho(X)$.

\section{The tests the code uses now}

\subsection{Everything grows along a path}

\begin{Lemma}
  \label{lem:grow}
  Let $A\ge0$ have no zero row, and $F=P+e_{ab}$.  Then
  $\nrm{FA}=\nrm{A}+s_b\ge\nrm{A}+1$, where $s_b$ is the sum of row~$b$
  of~$A$; and $c(FA)\ge c(A)$ and $r(FA)\ge r(A)$.
\end{Lemma}

\begin{proof}
  $FA=PA+e_{ab}A$, and $e_{ab}A$ is the matrix whose row~$a$ is row~$b$
  of~$A$, all else zero.  So the column sums of $FA$ are those of $A$ plus
  the entries of row~$b$ of $A$, and the row sums of $FA$ are those of $A$
  in a new order, one of them increased by~$s_b$.
\end{proof}

The matrix of a prefix dominates a permutation matrix, so it has no zero
row.  Hence along a path the norm and the smallest row and column sums
never decrease, and a prefix of length $L$ has norm at least $d+L$.  This
is what makes each of the next three tests safe: a bound that holds for
the closed path's matrix $BA$, and that $A$ already violates, stays
violated whatever the completion.

\subsection{The Ham--Song norm bound}

\begin{Lemma}[{Ham--Song~\cite[Lemma 3.1]{HamSong}}]
  \label{lem:hs}
  If $M$ is a $d\times d$ Perron--Frobenius matrix with
  $\lam=\rho(M)>1$, then $\nrm{M}\le\lam^d+d-1$.
\end{Lemma}

The proof takes a Perron eigenvector $v$ and its smallest coordinate
$v_p$.  Then $\lam^d v_p=(M^dv)_p\ge v_p\sum_j(M^d)_{pj}$, so $\lam^d$ is at
least the number of walks of length $d$ from $p$ in the multigraph of $M$.
A spanning tree at $p$ has $d-1$ edges, and every other edge is where some
such walk first leaves the tree.  The lemma improves an earlier bound of
Papadopoulos and Penner~\cite{PapadopoulosPenner}.

\emph{The test.}  Abandon a prefix when $\nrm{A}>\lfloor\Lam^d\rfloor+d-1$
(\code{find\_maxnorm}, \code{check\_all\_norms}).  Every closed path through
it then has $\nrm{BA}\ge\nrm{A}$ above the bound, so by Lemma~\ref{lem:hs}
its dilatation exceeds~$\Lam$.  Since the norm grows by at least one per
fold, the same number also caps the length of a path, and the code uses it
as \code{max\_path\_length}.  That cap is the norm test restated and never
cuts a path the norm test has not already cut.

\subsection{Row and column sums}

By~\eqref{eq:cw}, $r(A)$ and $c(A)$ are lower bounds for $\rho(A)$, and
along a path they never decrease; so abandoning a prefix when $r(A)>\Lam$
or $c(A)>\Lam$ is safe.  This is the test described by Ham and Song and in
the ttauto paper.

\subsection{The symmetrised norm}
\label{sec:symnorm}

With the macro \code{TTAUTO\_CHECK\_SYMMETRIC\_NORM}, the code used to
abandon a prefix also when $\frac1d\sum_{ij}\sqrt{A_{ij}A_{ji}}>\Lam$, with the
comment that this works only for Rauzy classes.  Here is why it is a lower
bound, and what the restriction amounts to.  Put
$G_{ij}=\sqrt{A_{ij}A_{ji}}$, the entrywise geometric mean of $A$ and
$A^{\mathsf T}$.  Then
\[
  \tfrac1d\textstyle\sum_{ij}G_{ij}
  = \tfrac1d\,\mathbf 1^{\mathsf T}G\,\mathbf 1 \le \rho(G)
  \le \rho(A)^{1/2}\rho(A^{\mathsf T})^{1/2} = \rho(A),
\]
the first inequality because $G$ is symmetric (the Rayleigh quotient), the
second by the inequality of Elsner, Johnson and Dias da Silva for weighted
geometric means of nonnegative matrices~\cite{EJD}, which generalises a
result of Karlin and Ost.  So it is a lower bound on $\rho(A)$ for
\emph{every} nonnegative matrix.  What needs Rauzy classes, meaning $P=I$
in~\eqref{eq:fold}, is the growth along a path: only then is $FA\ge A$, so
that $G$ never decreases.  In that case, though, $\rho(A)$ itself never
decreases (Proposition~\ref{prop:rhomono}), and it is the stronger test.
From four punctures on most folds do carry permutations, so there the test
could abandon a prefix that leads to an accepted path; it has therefore
been removed from the code (issue~\#23).

\subsection{Bad words}

This one is of a different kind.  Ham and Song's Lemma~3.3~\cite{HamSong}
says, roughly, that if repeating a loop $\gamma$ more times makes its
matrix larger without changing its pattern of zeros, then shortening the
repeat keeps a closed path pseudo-Anosov and does not raise its
dilatation.  So a search for the \emph{minimum} dilatation may skip paths
with long repeats: whatever they give, a shorter path gives no more.  It
is not safe in our sense, since those paths may well be accepted; it only
guarantees that the minimum is not among them.  A search that lists
everything in the window must not use it (issue~\#21).

\section{Where the tests lose}

\begin{Example}[Shears]
  \label{ex:shear}
  In dimension $d=2$ take $R=I+e_{12}$ and $L=I+e_{21}$, the two folds of
  the automaton for three punctures.  The prefix $R^k$ has
  $\nrm{R^k}=k+2$, $r(R^k)=c(R^k)=1$ and $\rho(R^k)=1$, so it passes every
  test above until $k+2>\Lam^2+1$.  But every completion that closes it
  into a primitive matrix contains an $L$, and then the trace is at least
  $k+2$ (Theorem~\ref{thm:d2}).  So nothing beyond $k\approx\Lam$ can
  succeed, while the search goes on to $k\approx\Lam^2$.
\end{Example}

\ttauto\ itself shows this on the three-puncture automaton.  These are its
runs to trace~$T$ from the benchmark of issue~\#22, all complete:

\begin{center}
  \small
  \begin{tabular}{rrrrr}
    \toprule
    $T$ & longest accepted & mean length & \code{max\_path\_length} &
      prefixes explored \\
    & path & of prefixes explored & & \\
    \midrule
    20 & 19 & 71 & 399 & 5\,230 \\
    40 & 39 & 214 & 1\,599 & 26\,430 \\
    80 & 79 & 677 & 6\,399 & 128\,478 \\
    \bottomrule
  \end{tabular}
\end{center}

\begin{Example}[Lemma~\ref{lem:hs} is not the culprit]
  $M=\begin{pmatrix}k&1\\k^2-1&k\end{pmatrix}$ is a product of $R$'s and
  $L$'s (every nonnegative $2\times2$ matrix of determinant one is), with
  trace $2k$, so $\lam\approx2k$ and $\nrm{M}=k^2+2k\approx\lam^2/4$.  So
  even among these matrices Lemma~\ref{lem:hs} is sharp up to a factor of
  about four.  The loss is in applying it to prefixes: the norm of a run
  of shears says almost nothing about the dilatation of its completions.
\end{Example}

\section{A bound over completions}

\subsection{Rauzy type: every $P$ is the identity}

\begin{Proposition}
  \label{prop:rhomono}
  If every fold is of the form $I+e_{ab}$, then $\rho(FA)\ge\rho(A)$.  So
  abandoning a prefix when $\rho(A)>\Lam$ is safe, and this test is at
  least as strong as the row sum, column sum and symmetrised norm tests.
\end{Proposition}

\begin{proof}
  $FA=A+e_{ab}A\ge A$.  Each of the other three is a lower bound
  for~$\rho(A)$.
\end{proof}

This does not help with Example~\ref{ex:shear}, since $\rho(R^k)=1$.  What
does help is that a completion has to \emph{mix}.

\begin{Lemma}
  \label{lem:expand}
  Let $B=(I+E_k)\cdots(I+E_1)$, each $E_i$ a single $e_{ab}$, and let $S$ be
  the set of their positions.  Then $B\ge I+E_S$, and $B_{ij}>0$ only if
  $i$ can be reached from $j$ in $\D(S)$ (every vertex reaching itself).
\end{Lemma}

\begin{proof}
  Expanding the product gives a sum of nonnegative terms, and those with at
  most one factor $E_i$ add up to $I+\sum_iE_i\ge I+E_S$.  A term
  $E_{i_r}\cdots E_{i_1}$ is nonzero only when the positions chain up,
  which is a walk in $\D(S)$.
\end{proof}

Call a set of positions $S$ \emph{admissible} for $A$ when
$\D(A)\cup\D(S)$ is strongly connected.

\begin{Theorem}
  \label{thm:rauzy}
  Let $A$ be the matrix of a prefix, and $\Pos$ a set of positions
  containing the position of every fold that can occur in a completion.
  Put
  \[
    \beta(A)=\min\bigl\{\rho\bigl((I+E_S)A\bigr) :
      S\subseteq\Pos \text{ admissible for } A\bigr\},
  \]
  with $\min\emptyset=\infty$.  Then $\rho(BA)\ge\beta(A)$ for every
  completion $B$ such that $BA$ is irreducible.  In particular abandoning
  the prefix when $\beta(A)>\Lam$ is safe.
\end{Theorem}

\begin{proof}
  Let $S$ be the positions of the folds in $B$.  If $(BA)_{ij}>0$ there is
  a $k$ with $A_{kj}>0$ and $B_{ik}>0$, so by Lemma~\ref{lem:expand} each
  edge $j\to i$ of $\D(BA)$ is a walk in $\D(A)\cup\D(S)$.  If $BA$ is
  irreducible, $\D(BA)$ is strongly connected, and hence so is
  $\D(A)\cup\D(S)$: $S$ is admissible.  And $BA\ge(I+E_S)A$ by
  Lemma~\ref{lem:expand}, so $\rho(BA)\ge\rho((I+E_S)A)\ge\beta(A)$.  An
  accepted matrix is primitive, hence irreducible, so none is lost.  The
  empty completion is the case $S=\emptyset$.
\end{proof}

Since $\rho((I+E_S)A)$ can only grow with $S$, the minimum is attained at a
minimal admissible~$S$.  In an automaton, $\Pos$ may depend on where the
prefix ends: it need only contain the positions that occur on paths from
there back to the start.

\begin{Theorem}[Dimension two]
  \label{thm:d2}
  A word $W$ in $R$ and $L$ containing both letters has
  $\tr W\ge|W|+1$, with equality for instance for $LR^m$.  So an accepted
  word, whose trace is at most $T=\Lam+\Lam^{-1}$, has length at most
  $T-1$.  Moreover $\beta(R^k)=\rho(LR^k)$: the test of
  Theorem~\ref{thm:rauzy} keeps the prefix $R^k$ exactly when $LR^k$ is
  itself accepted, so it cuts a run of shears at the length of the longest
  accepted word.
\end{Theorem}

\begin{proof}
  If $W$ contains both letters then $W\ge I+e_{12}+e_{21}$ by
  Lemma~\ref{lem:expand}, so $W_{12},W_{21}\ge1$.  Now $\tr(RW)=\tr
  W+W_{21}$ and $\tr(LW)=\tr W+W_{12}$, so each letter added on the left
  raises the trace by at least one.  The shortest right factor of $W$ that
  contains both letters is $LR^m$ or $RL^m$, whose trace is $m+2$, one more
  than its length.  For the last claim: $\D(R^k)$ has the single edge
  $2\to1$, so an admissible $S$ must contain $(2,1)$, the position of~$L$,
  and $(I+e_{21})R^k=LR^k$.
\end{proof}

\subsection{General folds}

When the folds carry permutations, move the permutations to the left.

\begin{Lemma}
  \label{lem:frame}
  Write a fold as $F=P(I+e_{a'b})$, where $P(a')=a$.  If a prefix has matrix
  $A=QN$, with $Q$ a permutation matrix and $N$ a product of matrices
  $I+e$, then
  \[
    FA=(PQ)\,\bigl(I+Q^{-1}e_{a'b}Q\bigr)N ,
  \]
  and $Q^{-1}e_{a'b}Q$ is again a single~$1$, at the position
  $(Q^{-1}(a'),Q^{-1}(b))$.  So, starting from the empty path, every path
  has matrix $QN$, where $Q$ is the product of its folds' permutation parts
  and $N$ is a product of factors $I+\tilde e$, one for each fold, with the
  units $\tilde e$ written in the labels of the start vertex.
\end{Lemma}

Apply this to a closed path that begins with the prefix, that is, the
prefix followed by a completion~$B$:
\[
  M = BA = Q'\,N'N .
\]
Here $Q'$ is the product of the permutation parts of \emph{all} the folds,
prefix and completion together (the path's \emph{total permutation}), and
$N'$ is the product of the factors $I+\tilde e$ contributed by the folds
of the completion.  Neither $Q'$ nor $N'$ is known while we are still at
the prefix, so the bound has to cover all the values they can take.

\begin{Theorem}
  \label{thm:general}
  Let the prefix have matrix $A=QN$ as in Lemma~\ref{lem:frame}.  Let
  $\Hgp$ be a set of permutations containing the total permutation of every
  closed path that begins with the prefix, and $\Pos$ a set of positions
  containing the position of every unit $\tilde e$ that a completion can
  contribute.  For $Q'\in\Hgp$ and $S\subseteq\Pos$, let $G_{Q',S}$ be the
  digraph with an edge $j\to Q'(i)$ whenever $j\to k$ in $\D(N)$ and $i$
  can be reached from $k$ in $\D(S)$.  Then for every completion $B$ such
  that the closed path's matrix $M=BA$ is irreducible,
  \[
    \rho(M)\ \ge\ \min\bigl\{\rho\bigl(Q'(I+E_S)N\bigr) :
      Q'\in\Hgp,\ S\subseteq\Pos,\ G_{Q',S}\ \text{strongly connected}\bigr\}.
  \]
  In particular abandoning the prefix when this minimum exceeds $\Lam$ is
  safe.
\end{Theorem}

\begin{proof}
  Write $M=Q'N'N$ as above, and let $S$ be the positions of the units
  in~$N'$.  By Lemma~\ref{lem:expand}, $N'\ge I+E_S$, so $M\ge Q'(I+E_S)N$
  and $\rho(M)\ge\rho(Q'(I+E_S)N)$.  Also by Lemma~\ref{lem:expand},
  $N'_{ik}>0$ only if $i$ can be reached from $k$ in $\D(S)$.  Now
  $M_{Q'(i)\,j}=(N'N)_{ij}$, so if it is positive there is a $k$ with
  $N_{kj}>0$ and $N'_{ik}>0$, which gives an edge $j\to Q'(i)$ of
  $G_{Q',S}$.  So $\D(M)$ is contained in $G_{Q',S}$, and if $M$ is
  irreducible $G_{Q',S}$ is strongly connected: the pair $(Q',S)$ is one of
  those in the minimum.
\end{proof}

With $\Hgp=\{I\}$ this is Theorem~\ref{thm:rauzy}.  Any sets $\Hgp$ and
$\Pos$ large enough will do; larger ones only weaken the bound.  In a
one-vertex automaton, where every word is a closed path, one can take for
$\Hgp$ the group generated by the permutation parts of the folds, and for
$\Pos$ the positions of the conjugates $q^{-1}e_{a'b}q$, $q\in\Hgp$, of
the folds' units.  In a general automaton the smallest choices are the
total permutations and units that occur on paths from where the prefix
ends back to the start; they can be computed once, on the automaton lifted
to pairs (vertex, permutation).

We write $\beta(A)$ for the minimum in Theorem~\ref{thm:general} too; with
$\Hgp=\{I\}$ it is the $\beta(A)$ of Theorem~\ref{thm:rauzy}.

\begin{Proposition}
  \label{prop:dominate}
  $\beta(A)\ge\max\bigl(r(A),c(A)\bigr)$, and if every fold is of the form
  $I+e_{ab}$ then $\beta(A)\ge\rho(A)$.  So the new test never keeps a
  prefix that the row sum or column sum test would abandon, nor, without
  permutations, one that the tests of Proposition~\ref{prop:rhomono} or
  the symmetrised norm would abandon.
\end{Proposition}

\begin{proof}
  Each candidate in the minimum is $\rho(Q'X)$ with $X=(I+E_S)N$.  A
  permutation $Q'$ leaves the column sums of $X$ unchanged and only
  reorders its row sums, so $c(Q'X)=c(X)$ and $r(Q'X)=r(X)$.  Since
  $X\ge N$, these are at least $c(N)$ and $r(N)$, which equal $c(A)$ and
  $r(A)$ because $A=QN$ is $N$ with its rows reordered.  By~\eqref{eq:cw},
  $\rho(Q'X)\ge\max(r(A),c(A))$.  Without permutations the candidates are
  $\rho((I+E_S)A)$ with $(I+E_S)A\ge A$, so each is at least $\rho(A)$,
  which is itself at least the symmetrised norm
  (Section~\ref{sec:symnorm}).
\end{proof}

\begin{Remark}
  \label{rem:norm}
  The Ham--Song norm test is not covered.  Lemma~\ref{lem:hs} applies to a
  Perron--Frobenius matrix, and the matrices $(I+E_S)A$ over which $\beta$
  is minimised need not even be irreducible: $S$ being admissible makes
  $\D(A)\cup\D(S)$ strongly connected, but an edge of $\D((I+E_S)A)$ is an
  edge of $\D(A)$ followed by at most one edge of $\D(S)$.  Whether a
  prefix can have a norm above the Ham--Song bound and yet
  $\beta(A)\le\Lam$ is open.  It matters little in practice: the norm test
  costs one sum and is what guarantees that the search ends, so it stays in
  any case.
\end{Remark}

\section{A check that knows nothing of \ttauto}
\label{sec:check}

\code{bounds\_toy.cpp} uses only the C++ standard library;
\code{bounds\_toy.py} is a slower version that accepts the same words and,
except for $C_1$ (see below), visits the same prefixes.  Each
takes a set of fold matrices as a one-vertex automaton, so that every word
is a closed path, and searches depth first, accepting the words whose
product is primitive with $1<\rho\le\Lam$.  Spectral radii are never just
estimated: they are bracketed by~\eqref{eq:cw}, iterating
$x\leftarrow(X+I)x$.  A decision to abandon a prefix uses only the lower
end of the bracket, so it is safe however few iterations are run; a
decision to accept uses the upper end, on a primitive matrix.

\emph{The reference list.}  For each set of folds and each $\Lam$ we first
run the search with the Ham--Song norm test alone.  By Lemma~\ref{lem:hs}
that test is safe, so this search finds every accepted word: its output is
the complete list, against which each other test is judged.

\emph{The tests compared}, each run as a separate search:
\begin{description}
\item[$C_0$] what \ttauto\ does now: the norm, row sum and column sum
  tests;
\item[$C_1$] $C_0$ together with $\rho(A)>\Lam$
  (Proposition~\ref{prop:rhomono}; used only where every $P$ is the
  identity);
\item[$C_2$] $C_0$ together with the new test of
  Theorem~\ref{thm:general}, taking the minimum over all subsets $S$ of
  $\Pos$ and all $Q'$ in $\Hgp$.  By Proposition~\ref{prop:dominate} the
  row and column sums never abandon a prefix that the new test keeps; they
  are kept because they are cheap and run first, so that the new test is
  computed only for prefixes they let through.  The norm test is kept
  because it guarantees that the search ends;
\item[faulty] $C_0$ together with $\rho(A)>\Lam$, used where the folds
  carry permutations.  Nothing above justifies it there.  We include it on
  purpose, as a test we expect to be unsafe, to show that the comparison
  would notice an unsafe test.  For this row, missing words is the
  expected outcome.
\end{description}

\emph{The columns of Tables~\ref{tab:RL}--\ref{tab:perm3}.}
\emph{Accepted} and \emph{classes} count the words found and their classes
up to cyclic rotation.  \emph{Visited} counts the prefixes the search
enters, and \emph{wasted} those among them that begin no accepted word:
the work a perfect test would have avoided.  \emph{Fewer than $C_0$} is
$C_0$'s count of visited prefixes divided by the test's.  \emph{Longest
prefix} is the longest prefix entered, and \emph{longest word} the longest
accepted word.  \emph{Check} says ``finds all'' when the test's accepted
words are exactly those of the reference list, and how many it misses
otherwise.

\emph{Shading.}  Green marks where the new test $C_2$ stands out: it
visits at least ten times fewer prefixes than $C_0$; it wastes none; or
its longest prefix is the longest accepted word while $C_0$ goes further.
Orange marks the faulty test missing words, which is the check working as
intended.

% Generated by make_tables.py; do not edit by hand.
\begin{table}[ht]
\centering
\small
\begin{tabular}{rlrrrrrrrl}
\toprule
& & & & & & fewer & longest & longest & \\
$\Lambda$ & test & accepted & classes & visited & wasted & than $C_0$ & prefix & word & check \\
\midrule
3 & norm only & 2 & 1 & 31 & 26 &  & 8 & 2 & reference \\
 & $C_0$ (current) & 2 & 1 & 27 & 22 &  & 8 &  & finds all \\
 & $C_1$ & 2 & 1 & 19 & 14 & 1.4$\times$ & 8 &  & finds all \\
 & $C_2$ (new) & 2 & 1 & 5 & \good{0} & 5.4$\times$ & \good{2} &  & finds all \\
\midrule
5 & norm only & 16 & 5 & 211 & 188 &  & 24 & 4 & reference \\
 & $C_0$ (current) & 16 & 5 & 103 & 80 &  & 24 &  & finds all \\
 & $C_1$ & 16 & 5 & 65 & 42 & 1.6$\times$ & 24 &  & finds all \\
 & $C_2$ (new) & 16 & 5 & 23 & \good{0} & 4.5$\times$ & \good{4} &  & finds all \\
\midrule
8 & norm only & 68 & 15 & 1\,307 & 1\,226 &  & 63 & 7 & reference \\
 & $C_0$ (current) & 68 & 15 & 321 & 240 &  & 63 &  & finds all \\
 & $C_1$ & 68 & 15 & 195 & 114 & 1.6$\times$ & 63 &  & finds all \\
 & $C_2$ (new) & 68 & 15 & 81 & \good{0} & 4.0$\times$ & \good{7} &  & finds all \\
\midrule
12 & norm only & 186 & 30 & 6\,441 & 6\,234 &  & 143 & 11 & reference \\
 & $C_0$ (current) & 186 & 30 & 837 & 630 &  & 143 &  & finds all \\
 & $C_1$ & 186 & 30 & 473 & 266 & 1.8$\times$ & 143 &  & finds all \\
 & $C_2$ (new) & 186 & 30 & 207 & \good{0} & 4.0$\times$ & \good{11} &  & finds all \\
\midrule
20 & norm only & 670 & 78 & 49\,077 & 48\,370 &  & 399 & 19 & reference \\
 & $C_0$ (current) & 670 & 78 & 2\,809 & 2\,102 &  & 399 &  & finds all \\
 & $C_1$ & 670 & 78 & 1\,469 & 762 & 1.9$\times$ & 399 &  & finds all \\
 & $C_2$ (new) & 670 & 78 & 707 & \good{0} & 4.0$\times$ & \good{19} &  & finds all \\
\bottomrule
\end{tabular}
\caption{$\{R,L\}$, $d=2$.  Each row is a separate search; the columns are explained in Section~\ref{sec:check}.  Shaded green: where the new test $C_2$ stands out.}
\label{tab:RL}
\end{table}

\begin{table}[ht]
\centering
\small
\begin{tabular}{rlrrrrrrrl}
\toprule
& & & & & & fewer & longest & longest & \\
$\Lambda$ & test & accepted & classes & visited & wasted & than $C_0$ & prefix & word & check \\
\midrule
2.2 & norm only & 0 & 0 & 48\,091 & 48\,091 &  & 9 & 0 & reference \\
 & $C_0$ (current) & 0 & 0 & 45\,001 & 45\,001 &  & 9 &  & finds all \\
 & $C_1$ & 0 & 0 & 40\,561 & 40\,561 & 1.1$\times$ & 9 &  & finds all \\
 & $C_2$ (new) & 0 & 0 & \good{1} & 1 & \good{45\,001$\times$} & 0 &  & finds all \\
\midrule
2.5 & norm only & 6 & 2 & 3\,006\,835 & 3\,006\,816 &  & 14 & 3 & reference \\
 & $C_0$ (current) & 6 & 2 & 2\,743\,123 & 2\,743\,104 &  & 14 &  & finds all \\
 & $C_1$ & 6 & 2 & 2\,540\,047 & 2\,540\,028 & 1.1$\times$ & 14 &  & finds all \\
 & $C_2$ (new) & 6 & 2 & \good{19} & \good{0} & \good{144\,375$\times$} & \good{3} &  & finds all \\
\bottomrule
\end{tabular}
\caption{all $I+e_{ab}$, $d=3$.  Each row is a separate search; the columns are explained in Section~\ref{sec:check}.  Shaded green: where the new test $C_2$ stands out.}
\label{tab:rauzy3}
\end{table}

\begin{table}[ht]
\centering
\small
\begin{tabular}{rlrrrrrrrl}
\toprule
& & & & & & fewer & longest & longest & \\
$\Lambda$ & test & accepted & classes & visited & wasted & than $C_0$ & prefix & word & check \\
\midrule
2.0 & norm only & 6 & 4 & 242 & 233 &  & 7 & 2 & reference \\
 & $C_0$ (current) & 6 & 4 & 236 & 227 &  & 7 &  & finds all \\
 & $C_2$ (new) & 6 & 4 & \good{16} & 7 & \good{14.8$\times$} & 3 &  & finds all \\
 & faulty, on purpose & 6 & 4 & 72 & 63 & 3.3$\times$ & 7 &  & finds all \\
\midrule
2.5 & norm only & 16 & 8 & 3\,607 & 3\,587 &  & 14 & 3 & reference \\
 & $C_0$ (current) & 16 & 8 & 2\,541 & 2\,521 &  & 14 &  & finds all \\
 & $C_2$ (new) & 16 & 8 & \good{55} & 35 & \good{46.2$\times$} & 5 &  & finds all \\
 & faulty, on purpose & 16 & 8 & 699 & 679 & 3.6$\times$ & 14 &  & finds all \\
\midrule
3.0 & norm only & 69 & 21 & 71\,582 & 71\,496 &  & 26 & 5 & reference \\
 & $C_0$ (current) & 69 & 21 & 47\,084 & 46\,998 &  & 26 &  & finds all \\
 & $C_2$ (new) & 69 & 21 & \good{173} & 87 & \good{272$\times$} & 8 &  & finds all \\
 & faulty, on purpose & 63 & 21 & 14\,804 & 14\,732 & 3.2$\times$ & 26 &  & \caught{misses 6} \\
\bottomrule
\end{tabular}
\caption{shears and a 3-cycle, $d=3$.  Each row is a separate search; the columns are explained in Section~\ref{sec:check}.  Shaded green: where the new test $C_2$ stands out.}
\label{tab:perm3}
\end{table}



\subsection*{What the tables show}

\emph{Nothing is missed.}  $C_0$, $C_1$ and $C_2$ find exactly the
reference list in every case.  The faulty test misses 6 of the 69 accepted
words at $\Lam=3$ in Table~\ref{tab:perm3}, so an unsafe test is caught.
At $\Lam=2$ and $2.5$ it happens to miss nothing, a reminder that one
passing case proves little.

\emph{Dimension two} (Table~\ref{tab:RL}).  $C_2$ wastes no prefix at any
$\Lam$, and its longest prefix is the longest accepted word, as
Theorem~\ref{thm:d2} predicts; $C_0$ goes about $\Lam^2$ deep.  At
$\Lam=20$, $C_2$ visits 707 prefixes, four times fewer than $C_0$ and 69
times fewer than the norm test alone.  The gain is modest here only because
$C_0$'s row and column sums already do much of the work in dimension two.
The classes found, 1, 5, 15, 30 and 78 up to traces 3, 5, 8, 12 and 20,
are those of the three-puncture benchmark of issue~\#22.

\emph{Dimension three, all $I+e_{ab}$} (Table~\ref{tab:rauzy3}).  This is
where the gain is dramatic.  At $\Lam=2.5$, $C_2$ visits 19 prefixes, none
wasted, where $C_0$ visits 2\,743\,123: about $1.4\times10^5$ times fewer.
At $\Lam=2.2$ nothing is accepted, and $C_2$ abandons every one-letter
prefix.  At $\Lam=3$, $C_2$ accepts 30 words (8 classes) after visiting 67
prefixes, 6 of them wasted, in a hundredth of a second.  $C_0$ did not
finish that case within 71 minutes, when it was stopped, and the reference
search did not finish in 100 seconds.  So it is not in the table, and
whether $C_2$ missed anything there is unchecked; but the contrast between
a hundredth of a second and more than an hour is itself the result.

\emph{With permutations} (Table~\ref{tab:perm3}).  $C_2$ now wastes some
prefixes, because it must allow for every possible total permutation, but
at $\Lam=3$ it still visits 173 prefixes where $C_0$ visits 47\,084, 272
times fewer.

\emph{$C_1$.}  Proposition~\ref{prop:rhomono} gains a factor of about two
in dimension two and little in dimension three: along a run of shears the
prefix's own dilatation stays near one.  The C++ bracket is also cautious
on reducible prefixes, where the Collatz--Wielandt iteration need not
reach $\rho$: at $\Lam=2.2$ the Python version, which computes the
eigenvalues, visits 36\,721 prefixes against 40\,561 here.  Both find the
complete list; the C++ one just abandons fewer hopeless prefixes.

\emph{Cost.}  $C_2$ as implemented tries every subset of $\Pos$ and every
element of $\Hgp$ at every prefix.  Here $|\Pos|\le6$ and $|\Hgp|\le3$, so
this is cheap.  In \ttauto's automata $d$ reaches~8, so $\Pos$ can have up
to $d(d-1)=56$ positions and $\Hgp$ can be large.

\section{Open questions}
\label{sec:open}

\begin{enumerate}
\item \emph{The cost of the new test.}  Only minimal admissible sets
  matter.  When $\D(A)$ is nearly strongly connected there are few; when it
  is sparse there are many, but then the bound tends to be small and to
  abandon little, so it could be skipped.  Better: is there a single
  matrix, computable in polynomial time, lying below $(I+E_S)A$ for every
  admissible $S$?  The strongly connected components of $\D(A)$ suggest
  one, since every admissible $S$ must cross each cut between them.
\item \emph{Permutations in a real automaton.}  $\Hgp$ and $\Pos$ come from
  the automaton lifted to (vertex, permutation) states.  How large are they
  for \ttauto's strata?  A cheaper bound, $\min_{Q'\in\Hgp}\rho(Q'N)$,
  needs no irreducibility and already extends $C_1$ to folds with
  permutations, since $M=Q'N'N\ge Q'N$.
\item \emph{The norm test.}  Does $\beta(A)>\Lam$ whenever the norm of
  $A$ exceeds the Ham--Song bound (Remark~\ref{rem:norm})?  A positive answer would make the new
  test a single statement containing all the others.
\item \emph{A length cap.}  Theorem~\ref{thm:d2} caps an accepted word at
  $T-1$ in dimension two.  No cap logarithmic in $\Lam$ is possible in
  general, since a run of shears grows only linearly.  Is the longest
  accepted path polynomial in $\Lam$ at fixed $d$, and with what exponent?
  The program reports the longest accepted word, a place to start.
\item \emph{Beyond \ttauto.}  Rauzy--Veech induction matrices are exactly
  of the form $I+e_{ab}$, so Theorem~\ref{thm:rauzy} applies as it stands
  to searches over Rauzy classes, such as minimal dilatations in strata of
  abelian differentials.  Theorem~\ref{thm:general} applies to any folding
  automaton, Ham and Song's included.
\item \emph{Into \ttauto.}  Measure $\Hgp$ and $\Pos$ on the real automata,
  implement $C_2$ or a cheaper form of it in \code{check\_all\_norms}, and
  require that it reproduce the lists known to be complete: the 2\,541
  classes of the three-puncture benchmark, the 10 and 6 classes of
  issue~\#16, and complete windows at four punctures.
\end{enumerate}

\begin{thebibliography}{9}

\bibitem{EJD}
L.~Elsner, C.~R.~Johnson and J.~A.~Dias da Silva, \emph{The Perron root of
  a weighted geometric mean of nonnegative matrices}, Linear Multilinear
  Algebra \textbf{24} (1988), 1--13.

\bibitem{HamSong}
J.-Y.~Ham and W.~T.~Song, \emph{The minimum dilatation of pseudo-Anosov
  5-braids}, Experiment.\ Math.\ \textbf{16} (2007), 167--179.

\bibitem{PapadopoulosPenner}
A.~Papadopoulos and R.~C.~Penner, \emph{Enumerating pseudo-Anosov
  foliations}, Pacific J.\ Math.\ \textbf{142} (1990), 159--173.

\bibitem{Seneta}
E.~Seneta, \emph{Non-negative matrices and Markov chains}, 2nd ed.,
  Springer, 1981.

\end{thebibliography}

\end{document}
